
        You are a research assistant AI tasked with generating a scientific paper based on provided literature. Follow these steps:
    1. Analyze the given References.
    2. Identify gaps in existing research to establish the motivation for a new study.
    3. Propose a main idea for a new research work.
    4. Write the paper's main content in LaTeX format, including all the following parts, please must have tables:
    - Title
    - Abstract
    - Introduction
    - Related Work
    - Methods/Experiment
    - Results with tables
    - Discussion
    - Conclusion
    - Contributions
    Ensure all content is original, academically rigorous, and follows standard scientific writing conventions.Please make all decision by yourself,do not ask for any instructions

        Here are the references to be used for generating the paper.
        You MUST base the paper on these references and integrate them naturally.

        References:
        --------------------
        @misc{mansour2026doinglessdataaugmentation,
      title={Doing More with Less: Data Augmentation for Sudanese Dialect Automatic Speech Recognition}, 
      author={Ayman Mansour},
      year={2026},
      eprint={2601.06802},
      archivePrefix={arXiv},
      primaryClass={cs.CL},
      url={https://arxiv.org/abs/2601.06802},
      abstract = {Although many Automatic Speech Recognition (ASR) systems have been developed for Modern Standard Arabic (MSA) and Dialectal Arabic (DA), few studies have focused on dialect-specific implementations, particularly for low-resource Arabic dialects such as Sudanese. This paper presents a comprehensive study of data augmentation techniques for fine-tuning OpenAI Whisper models and establishes the first benchmark for the Sudanese dialect. Two augmentation strategies are investigated: (1) self-training with pseudo-labels generated from unlabeled speech, and (2) TTS-based augmentation using synthetic speech from the Klaam TTS system. The best-performing model, Whisper-Medium fine-tuned with combined self-training and TTS augmentation (28.4 hours), achieves a Word Error Rate (WER) of 57.1\% on the evaluation set and 51.6\% on an out-of-domain holdout set substantially outperforming zero-shot multilingual Whisper (78.8\% WER) and MSA-specialized Arabic models (73.8-123\% WER). All experiments used low-cost resources (Kaggle free tier and Lightning.ai trial), demonstrating that strategic data augmentation can overcome resource limitations for low-resource dialects and provide a practical roadmap for developing ASR systems for low-resource Arabic dialects and other marginalized language varieties. The models, evaluation benchmarks, and reproducible training pipelines are publicly released to facilitate future research on low-resource Arabic ASR.}
}@misc{zheng2026mosaicunlockinglongcontextinference,
      title={Mosaic: Unlocking Long-Context Inference for Diffusion LLMs via Global Memory Planning and Dynamic Peak Taming}, 
      author={Liang Zheng and Bowen Shi and Yitao Hu and Jiawei Zhang and Ruofan Li and Sheng Chen and Wenxin Li and Keqiu Li},
      year={2026},
      eprint={2601.06562},
      archivePrefix={arXiv},
      primaryClass={cs.LG},
      url={https://arxiv.org/abs/2601.06562},
      abstract = {Diffusion-based large language models (dLLMs) have emerged as a promising paradigm, utilizing simultaneous denoising to enable global planning and iterative refinement. While these capabilities are particularly advantageous for long-context generation, deploying such models faces a prohibitive memory capacity barrier stemming from severe system inefficiencies. We identify that existing inference systems are ill-suited for this paradigm: unlike autoregressive models constrained by the cumulative KV-cache, dLLMs are bottlenecked by transient activations recomputed at every step. Furthermore, general-purpose memory reuse mechanisms lack the global visibility to adapt to dLLMs' dynamic memory peaks, which toggle between logits and FFNs. To address these mismatches, we propose Mosaic, a memory-efficient inference system that shifts from local, static management to a global, dynamic paradigm. Mosaic integrates a mask-only logits kernel to eliminate redundancy, a lazy chunking optimizer driven by an online heuristic search to adaptively mitigate dynamic peaks, and a global memory manager to resolve fragmentation via virtual addressing. Extensive evaluations demonstrate that Mosaic achieves an average 2.71$\\times$ reduction in the memory peak-to-average ratio and increases the maximum inference sequence length supportable on identical hardware by 15.89-32.98$\\times$. This scalability is achieved without compromising accuracy and speed, and in fact reducing latency by 4.12\%-23.26\%.}
}@misc{yang2026stephanie2thinkingwaitingmaking,
      title={Stephanie2: Thinking, Waiting, and Making Decisions Like Humans in Step-by-Step AI Social Chat}, 
      author={Hao Yang and Hongyuan Lu and Dingkang Yang and Wenliang Yang and Peng Sun and Xiaochuan Zhang and Jun Xiao and Kefan He and Wai Lam and Yang Liu and Xinhua Zeng},
      year={2026},
      eprint={2601.05657},
      archivePrefix={arXiv},
      primaryClass={cs.CL},
      url={https://arxiv.org/abs/2601.05657},
      abstract = {Instant-messaging human social chat typically progresses through a sequence of short messages. Existing step-by-step AI chatting systems typically split a one-shot generation into multiple messages and send them sequentially, but they lack an active waiting mechanism and exhibit unnatural message pacing. In order to address these issues, we propose Stephanie2, a novel next-generation step-wise decision-making dialogue agent. With active waiting and message-pace adaptation, Stephanie2 explicitly decides at each step whether to send or wait, and models latency as the sum of thinking time and typing time to achieve more natural pacing. We further introduce a time-window-based dual-agent dialogue system to generate pseudo dialogue histories for human and automatic evaluations. Experiments show that Stephanie2 clearly outperforms Stephanie1 on metrics such as naturalness and engagement, and achieves a higher pass rate on human evaluation with the role identification Turing test.}
}@misc{vamshi2026manifoldbasedsamplingincontexthallucination,
      title={Manifold-based Sampling for In-Context Hallucination Detection in Large Language Models}, 
      author={Bodla Krishna Vamshi and Rohan Bhatnagar and Haizhao Yang},
      year={2026},
      eprint={2601.06196},
      archivePrefix={arXiv},
      primaryClass={cs.LG},
      url={https://arxiv.org/abs/2601.06196},
      abstract = {Large language models (LLMs) frequently generate factually incorrect or unsupported content, commonly referred to as hallucinations. Prior work has explored decoding strategies, retrieval augmentation, and supervised fine-tuning for hallucination detection, while recent studies show that in-context learning (ICL) can substantially influence factual reliability. However, existing ICL demonstration selection methods often rely on surface-level similarity heuristics and exhibit limited robustness across tasks and models.   We propose MB-ICL, a manifold-based demonstration sampling framework for selecting in-context demonstrations that leverages latent representations extracted from frozen LLMs. By jointly modeling local manifold structure and class-aware prototype geometry, MB-ICL selects demonstrations based on their proximity to learned prototypes rather than lexical or embedding similarity alone.   Across factual verification (FEVER) and hallucination detection (HaluEval) benchmarks, MB-ICL outperforms standard ICL selection baselines in the majority of evaluated settings, with particularly strong gains on dialogue and summarization tasks. The method remains robust under temperature perturbations and model variation, indicating improved stability compared to heuristic retrieval strategies. While lexical retrieval can remain competitive in certain question-answering regimes, our results demonstrate that manifold-based prototype selection provides a reliable and training light approach for hallucination detection without modifying LLM parameters, offering a principled direction for improved ICL demonstration selection.}
}@misc{kennedy2026leftrightcenterevaluating,
      title={Left, Right, or Center? Evaluating LLM Framing in News Classification and Generation}, 
      author={Molly Kennedy and Ali Parker and Yihong Liu and Hinrich Schütze},
      year={2026},
      eprint={2601.05835},
      archivePrefix={arXiv},
      primaryClass={cs.CL},
      url={https://arxiv.org/abs/2601.05835},
      abstract = {Large Language Model (LLM) based summarization and text generation are increasingly used for producing and rewriting text, raising concerns about political framing in journalism where subtle wording choices can shape interpretation. Across nine state-of-the-art LLMs, we study political framing by testing whether LLMs' classification-based bias signals align with framing behavior in their generated summaries. We first compare few-shot ideology predictions against LEFT/CENTER/RIGHT labels. We then generate "steered" summaries under FAITHFUL, CENTRIST, LEFT, and RIGHT prompts, and score all outputs using a single fixed ideology evaluator. We find pervasive ideological center-collapse in both article-level ratings and generated text, indicating a systematic tendency toward centrist framing. Among evaluated models, Grok 4 is by far the most ideologically expressive generator, while Claude Sonnet 4.5 and Llama 3.1 achieve the strongest bias-rating performance among commercial and open-weight models, respectively.}
}@misc{tian2026stagebenchmarkknowledgegraph,
      title={STAGE: A Benchmark for Knowledge Graph Construction, Question Answering, and In-Script Role-Playing over Movie Screenplays}, 
      author={Qiuyu Tian and Yiding Li and Fengyi Chen and Zequn Liu and Youyong Kong and Fan Guo and Yuyao Li and Jinjing Shen and Zhijing Xie and Yiyun Luo and Xin Zhang},
      year={2026},
      eprint={2601.08510},
      archivePrefix={arXiv},
      primaryClass={cs.CL},
      url={https://arxiv.org/abs/2601.08510},
      abstract = {Movie screenplays are rich long-form narratives that interleave complex character relationships, temporally ordered events, and dialogue-driven interactions. While prior benchmarks target individual subtasks such as question answering or dialogue generation, they rarely evaluate whether models can construct a coherent story world and use it consistently across multiple forms of reasoning and generation. We introduce STAGE (Screenplay Text, Agents, Graphs and Evaluation), a unified benchmark for narrative understanding over full-length movie screenplays. STAGE defines four tasks: knowledge graph construction, scene-level event summarization, long-context screenplay question answering, and in-script character role-playing, all grounded in a shared narrative world representation. The benchmark provides cleaned scripts, curated knowledge graphs, and event- and character-centric annotations for 150 films across English and Chinese, enabling holistic evaluation of models' abilities to build world representations, abstract and verify narrative events, reason over long narratives, and generate character-consistent responses.}
}@misc{huang2026probingmultimodallargelanguage,
      title={Probing Multimodal Large Language Models on Cognitive Biases in Chinese Short-Video Misinformation}, 
      author={Jen-tse Huang and Chang Chen and Shiyang Lai and Wenxuan Wang and Michelle R. Kaufman and Mark Dredze},
      year={2026},
      eprint={2601.06600},
      archivePrefix={arXiv},
      primaryClass={cs.CL},
      url={https://arxiv.org/abs/2601.06600},
      abstract = {Short-video platforms have become major channels for misinformation, where deceptive claims frequently leverage visual experiments and social cues. While Multimodal Large Language Models (MLLMs) have demonstrated impressive reasoning capabilities, their robustness against misinformation entangled with cognitive biases remains under-explored. In this paper, we introduce a comprehensive evaluation framework using a high-quality, manually annotated dataset of 200 short videos spanning four health domains. This dataset provides fine-grained annotations for three deceptive patterns, experimental errors, logical fallacies, and fabricated claims, each verified by evidence such as national standards and academic literature. We evaluate eight frontier MLLMs across five modality settings. Experimental results demonstrate that Gemini-2.5-Pro achieves the highest performance in the multimodal setting with a belief score of 71.5/100, while o3 performs the worst at 35.2. Furthermore, we investigate social cues that induce false beliefs in videos and find that models are susceptible to biases like authoritative channel IDs.}
}@misc{kaneko2026paraphrasingadversarialattackllmasareviewer,
      title={Paraphrasing Adversarial Attack on LLM-as-a-Reviewer}, 
      author={Masahiro Kaneko},
      year={2026},
      eprint={2601.06884},
      archivePrefix={arXiv},
      primaryClass={cs.CL},
      url={https://arxiv.org/abs/2601.06884},
      abstract = {The use of large language models (LLMs) in peer review systems has attracted growing attention, making it essential to examine their potential vulnerabilities. Prior attacks rely on prompt injection, which alters manuscript content and conflates injection susceptibility with evaluation robustness. We propose the Paraphrasing Adversarial Attack (PAA), a black-box optimization method that searches for paraphrased sequences yielding higher review scores while preserving semantic equivalence and linguistic naturalness. PAA leverages in-context learning, using previous paraphrases and their scores to guide candidate generation. Experiments across five ML and NLP conferences with three LLM reviewers and five attacking models show that PAA consistently increases review scores without changing the paper's claims. Human evaluation confirms that generated paraphrases maintain meaning and naturalness. We also find that attacked papers exhibit increased perplexity in reviews, offering a potential detection signal, and that paraphrasing submissions can partially mitigate attacks.}
}@misc{wang2026orthogeolorageometricparameterefficientfinetuning,
      title={OrthoGeoLoRA: Geometric Parameter-Efficient Fine-Tuning for Structured Social Science Concept Retrieval on theWeb}, 
      author={Zeqiang Wang and Xinyue Wu and Chenxi Li and Zixi Chen and Nishanth Sastry and Jon Johnson and Suparna De},
      year={2026},
      eprint={2601.09185},
      archivePrefix={arXiv},
      primaryClass={cs.CL},
      url={https://arxiv.org/abs/2601.09185},
      abstract = {Large language models and text encoders increasingly power web-based information systems in the social sciences, including digital libraries, data catalogues, and search interfaces used by researchers, policymakers, and civil society. Full fine-tuning is often computationally and energy intensive, which can be prohibitive for smaller institutions and non-profit organizations in the Web4Good ecosystem. Parameter-Efficient Fine-Tuning (PEFT), especially Low-Rank Adaptation (LoRA), reduces this cost by updating only a small number of parameters. We show that the standard LoRA update $ΔW = BA^\\top$ has geometric drawbacks: gauge freedom, scale ambiguity, and a tendency toward rank collapse. We introduce OrthoGeoLoRA, which enforces an SVD-like form $ΔW = BΣA^\\top$ by constraining the low-rank factors to be orthogonal (Stiefel manifold). A geometric reparameterization implements this constraint while remaining compatible with standard optimizers such as Adam and existing fine-tuning pipelines. We also propose a benchmark for hierarchical concept retrieval over the European Language Social Science Thesaurus (ELSST), widely used to organize social science resources in digital repositories. Experiments with a multilingual sentence encoder show that OrthoGeoLoRA outperforms standard LoRA and several strong PEFT variants on ranking metrics under the same low-rank budget, offering a more compute- and parameter-efficient path to adapt foundation models in resource-constrained settings.}
}@misc{toraman2026turkbenchbenchmarkevaluatingturkish,
      title={TurkBench: A Benchmark for Evaluating Turkish Large Language Models}, 
      author={Çağrı Toraman and Ahmet Kaan Sever and Ayse Aysu Cengiz and Elif Ecem Arslan and Görkem Sevinç and Mete Mert Birdal and Yusuf Faruk Güldemir and Ali Buğra Kanburoğlu and Sezen Felekoğlu and Osman Gürlek and Sarp Kantar and Birsen Şahin Kütük and Büşra Tufan and Elif Genç and Serkan Coşkun and Gupse Ekin Demir and Muhammed Emin Arayıcı and Olgun Dursun and Onur Gungor and Susan Üsküdarlı and Abdullah Topraksoy and Esra Darıcı},
      year={2026},
      eprint={2601.07020},
      archivePrefix={arXiv},
      primaryClass={cs.CL},
      url={https://arxiv.org/abs/2601.07020},
      abstract = {With the recent surge in the development of large language models, the need for comprehensive and language-specific evaluation benchmarks has become critical. While significant progress has been made in evaluating English language models, benchmarks for other languages, particularly those with unique linguistic characteristics such as Turkish, remain less developed. Our study introduces TurkBench, a comprehensive benchmark designed to assess the capabilities of generative large language models in the Turkish language. TurkBench involves 8,151 data samples across 21 distinct subtasks. These are organized under six main categories of evaluation: Knowledge, Language Understanding, Reasoning, Content Moderation, Turkish Grammar and Vocabulary, and Instruction Following. The diverse range of tasks and the culturally relevant data would provide researchers and developers with a valuable tool for evaluating their models and identifying areas for improvement. We further publish our benchmark for online submissions at https://huggingface.co/turkbench}
}@misc{nehrdich2026mitrasamgrahacomprehensiveclassicalsanskrit,
      title={Mitrasamgraha: A Comprehensive Classical Sanskrit Machine Translation Dataset}, 
      author={Sebastian Nehrdich and David Allport and Sven Sellmer and Jivnesh Sandhan and Manoj Balaji Jagadeeshan and Pawan Goyal and Sujeet Kumar and Kurt Keutzer},
      year={2026},
      eprint={2601.07314},
      archivePrefix={arXiv},
      primaryClass={cs.CL},
      url={https://arxiv.org/abs/2601.07314},
      abstract = {While machine translation is regarded as a "solved problem" for many high-resource languages, close analysis quickly reveals that this is not the case for content that shows challenges such as poetic language, philosophical concepts, multi-layered metaphorical expressions, and more. Sanskrit literature is a prime example of this, as it combines a large number of such challenges in addition to inherent linguistic features like sandhi, compounding, and heavy morphology, which further complicate NLP downstream tasks. It spans multiple millennia of text production time as well as a large breadth of different domains, ranging from ritual formulas via epic narratives, philosophical treatises, poetic verses up to scientific material. As of now, there is a strong lack of publicly available resources that cover these different domains and temporal layers of Sanskrit. We therefore introduce Mitrasamgraha, a high-quality Sanskrit-to-English machine translation dataset consisting of 391,548 bitext pairs, more than four times larger than the largest previously available Sanskrit dataset Itih=asa. It covers a time period of more than three millennia and a broad range of historical Sanskrit domains. In contrast to web-crawled datasets, the temporal and domain annotation of this dataset enables fine-grained study of domain and time period effects on MT performance. We also release a validation set consisting of 5,587 and a test set consisting of 5,552 post-corrected bitext pairs. We conduct experiments benchmarking commercial and open models on this dataset and fine-tune NLLB and Gemma models on the dataset, showing significant improvements, while still recognizing significant challenges in the translation of complex compounds, philosophical concepts, and multi-layered metaphors. We also analyze how in-context learning on this dataset impacts the performance of commercial models}
}@misc{ma2026slamllmmodularopensourcemultimodal,
      title={SLAM-LLM: A Modular, Open-Source Multimodal Large Language Model Framework and Best Practice for Speech, Language, Audio and Music Processing}, 
      author={Ziyang Ma and Guanrou Yang and Wenxi Chen and Zhifu Gao and Yexing Du and Xiquan Li and Zhisheng Zheng and Haina Zhu and Jianheng Zhuo and Zheshu Song and Ruiyang Xu and Tiranrui Wang and Yifan Yang and Yanqiao Zhu and Zhikang Niu and Liumeng Xue and Yinghao Ma and Ruibin Yuan and Shiliang Zhang and Kai Yu and Eng Siong Chng and Xie Chen},
      year={2026},
      eprint={2601.09385},
      archivePrefix={arXiv},
      primaryClass={cs.SD},
      doi={https://doi.org/10.1109/JSTSP.2026.3653157},
      url={https://arxiv.org/abs/2601.09385},
      abstract = {The recent surge in open-source Multimodal Large Language Models (MLLM) frameworks, such as LLaVA, provides a convenient kickoff for artificial intelligence developers and researchers. However, most of the MLLM frameworks take vision as the main input modality, and provide limited in-depth support for the modality of speech, audio, and music. This situation hinders the development of audio-language models, and forces researchers to spend a lot of effort on code writing and hyperparameter tuning. We present SLAM-LLM, an open-source deep learning framework designed to train customized MLLMs, focused on speech, language, audio, and music processing. SLAM-LLM provides a modular configuration of different encoders, projectors, LLMs, and parameter-efficient fine-tuning plugins. SLAM-LLM also includes detailed training and inference recipes for mainstream tasks, along with high-performance checkpoints like LLM-based Automatic Speech Recognition (ASR), Automated Audio Captioning (AAC), and Music Captioning (MC). Some of these recipes have already reached or are nearing state-of-the-art performance, and some relevant techniques have also been accepted by academic papers. We hope SLAM-LLM will accelerate iteration, development, data engineering, and model training for researchers. We are committed to continually pushing forward audio-based MLLMs through this open-source framework, and call on the community to contribute to the LLM-based speech, audio and music processing.}
}@misc{xie2026roapreadingorderattentionpriorpipeline,
      title={ROAP: A Reading-Order and Attention-Prior Pipeline for Optimizing Layout Transformers in Key Information Extraction}, 
      author={Tingwei Xie and Jinxin He and Yonghong Song},
      year={2026},
      eprint={2601.05470},
      archivePrefix={arXiv},
      primaryClass={cs.CV},
      url={https://arxiv.org/abs/2601.05470},
      abstract = {The efficacy of Multimodal Transformers in visually-rich document understanding (VrDU) is critically constrained by two inherent limitations: the lack of explicit modeling for logical reading order and the interference of visual tokens that dilutes attention on textual semantics.   To address these challenges, this paper presents ROAP, a lightweight and architecture-agnostic pipeline designed to optimize attention distributions in Layout Transformers without altering their pre-trained backbones.   The proposed pipeline first employs an Adaptive-XY-Gap (AXG-Tree) to robustly extract hierarchical reading sequences from complex layouts. These sequences are then integrated into the attention mechanism via a Reading-Order-Aware Relative Position Bias (RO-RPB). Furthermore, a Textual-Token Sub-block Attention Prior (TT-Prior) is introduced to adaptively suppress visual noise and enhance fine-grained text-text interactions.   Extensive experiments on the FUNSD and CORD benchmarks demonstrate that ROAP consistently improves the performance of representative backbones, including LayoutLMv3 and GeoLayoutLM.   These findings confirm that explicitly modeling reading logic and regulating modality interference are critical for robust document understanding, offering a scalable solution for complex layout analysis. The implementation code will be released at https://github.com/KevinYuLei/ROAP.}
}@misc{sun2026cumaaligningllmssparse,
      title={CuMA: Aligning LLMs with Sparse Cultural Values via Demographic-Aware Mixture of Adapters}, 
      author={Ao Sun and Xiaoyu Wang and Zhe Tan and Yu Li and Jiachen Zhu and Shu Su and Yuheng Jia},
      year={2026},
      eprint={2601.04885},
      archivePrefix={arXiv},
      primaryClass={cs.CL},
      url={https://arxiv.org/abs/2601.04885},
      abstract = {As Large Language Models (LLMs) serve a global audience, alignment must transition from enforcing universal consensus to respecting cultural pluralism. We demonstrate that dense models, when forced to fit conflicting value distributions, suffer from \\textbf\{Mean Collapse\}, converging to a generic average that fails to represent diverse groups. We attribute this to \\textbf\{Cultural Sparsity\}, where gradient interference prevents dense parameters from spanning distinct cultural modes. To resolve this, we propose \\textbf\{\\textsc\{CuMA\}\} (\\textbf\{Cu\}ltural \\textbf\{M\}ixture of \\textbf\{A\}dapters), a framework that frames alignment as a \\textbf\{conditional capacity separation\} problem. By incorporating demographic-aware routing, \\textsc\{CuMA\} internalizes a \\textit\{Latent Cultural Topology\} to explicitly disentangle conflicting gradients into specialized expert subspaces. Extensive evaluations on WorldValuesBench, Community Alignment, and PRISM demonstrate that \\textsc\{CuMA\} achieves state-of-the-art performance, significantly outperforming both dense baselines and semantic-only MoEs. Crucially, our analysis confirms that \\textsc\{CuMA\} effectively mitigates mean collapse, preserving cultural diversity. Our code is available at https://github.com/Throll/CuMA.}
}@misc{cheng2026culturalcompassframeworkorganizing,
      title={Cultural Compass: A Framework for Organizing Societal Norms to Detect Violations in Human-AI Conversations}, 
      author={Myra Cheng and Vinodkumar Prabhakaran and Alice Oh and Hayk Stepanyan and Aishwarya Verma and Charu Kalia and Erin MacMurray van Liemt and Sunipa Dev},
      year={2026},
      eprint={2601.07973},
      archivePrefix={arXiv},
      primaryClass={cs.CY},
      url={https://arxiv.org/abs/2601.07973},
      abstract = {Generative AI models ought to be useful and safe across cross-cultural contexts. One critical step toward this goal is understanding how AI models adhere to sociocultural norms. While this challenge has gained attention in NLP, existing work lacks both nuance and coverage in understanding and evaluating models' norm adherence. We address these gaps by introducing a taxonomy of norms that clarifies their contexts (e.g., distinguishing between human-human norms that models should recognize and human-AI interactional norms that apply to the human-AI interaction itself), specifications (e.g., relevant domains), and mechanisms (e.g., modes of enforcement). We demonstrate how our taxonomy can be operationalized to automatically evaluate models' norm adherence in naturalistic, open-ended settings. Our exploratory analyses suggest that state-of-the-art models frequently violate norms, though violation rates vary by model, interactional context, and country. We further show that violation rates also vary by prompt intent and situational framing. Our taxonomy and demonstrative evaluation pipeline enable nuanced, context-sensitive evaluation of cultural norm adherence in realistic settings.}
}@misc{ahn2026enrichingsemanticprofilesknowledge,
      title={Enriching Semantic Profiles into Knowledge Graph for Recommender Systems Using Large Language Models}, 
      author={Seokho Ahn and Sungbok Shin and Young-Duk Seo},
      year={2026},
      eprint={2601.08148},
      archivePrefix={arXiv},
      primaryClass={cs.IR},
      url={https://arxiv.org/abs/2601.08148},
      abstract = {Rich and informative profiling to capture user preferences is essential for improving recommendation quality. However, there is still no consensus on how best to construct and utilize such profiles. To address this, we revisit recent profiling-based approaches in recommender systems along four dimensions: 1) knowledge base, 2) preference indicator, 3) impact range, and 4) subject. We argue that large language models (LLMs) are effective at extracting compressed rationales from diverse knowledge sources, while knowledge graphs (KGs) are better suited for propagating these profiles to extend their reach. Building on this insight, we propose a new recommendation model, called SPiKE. SPiKE consists of three core components: i) Entity profile generation, which uses LLMs to generate semantic profiles for all KG entities; ii) Profile-aware KG aggregation, which integrates these profiles into the KG; and iii) Pairwise profile preference matching, which aligns LLM- and KG-based representations during training. In experiments, we demonstrate that SPiKE consistently outperforms state-of-the-art KG- and LLM-based recommenders in real-world settings.}
}@misc{zhang2026miprunoptimizelargelanguage,
      title={MI-PRUN: Optimize Large Language Model Pruning via Mutual Information}, 
      author={Hao Zhang and Zhibin Zhang and Guangxin Wu and He Chen and Jiafeng Guo and Xueqi Cheng},
      year={2026},
      eprint={2601.07212},
      archivePrefix={arXiv},
      primaryClass={cs.CL},
      url={https://arxiv.org/abs/2601.07212},
      abstract = {Large Language Models (LLMs) have become indispensable across various domains, but this comes at the cost of substantial computational and memory resources. Model pruning addresses this by removing redundant components from models. In particular, block pruning can achieve significant compression and inference acceleration. However, existing block pruning methods are often unstable and struggle to attain globally optimal solutions. In this paper, we propose a mutual information based pruning method MI-PRUN for LLMs. Specifically, we leverages mutual information to identify redundant blocks by evaluating transitions in hidden states. Additionally, we incorporate the Data Processing Inequality (DPI) to reveal the relationship between the importance of entire contiguous blocks and that of individual blocks. Moreover, we develop the Fast-Block-Select algorithm, which iteratively updates block combinations to achieve a globally optimal solution while significantly improving the efficiency. Extensive experiments across various models and datasets demonstrate the stability and effectiveness of our method.}
}@misc{zhou2026detectingllmgeneratedtextperformance,
      title={Detecting LLM-Generated Text with Performance Guarantees}, 
      author={Hongyi Zhou and Jin Zhu and Ying Yang and Chengchun Shi},
      year={2026},
      eprint={2601.06586},
      archivePrefix={arXiv},
      primaryClass={cs.CL},
      url={https://arxiv.org/abs/2601.06586},
      abstract = {Large language models (LLMs) such as GPT, Claude, Gemini, and Grok have been deeply integrated into our daily life. They now support a wide range of tasks -- from dialogue and email drafting to assisting with teaching and coding, serving as search engines, and much more. However, their ability to produce highly human-like text raises serious concerns, including the spread of fake news, the generation of misleading governmental reports, and academic misconduct. To address this practical problem, we train a classifier to determine whether a piece of text is authored by an LLM or a human. Our detector is deployed on an online CPU-based platform https://huggingface.co/spaces/stats-powered-ai/StatDetectLLM, and contains three novelties over existing detectors: (i) it does not rely on auxiliary information, such as watermarks or knowledge of the specific LLM used to generate the text; (ii) it more effectively distinguishes between human- and LLM-authored text; and (iii) it enables statistical inference, which is largely absent in the current literature. Empirically, our classifier achieves higher classification accuracy compared to existing detectors, while maintaining type-I error control, high statistical power, and computational efficiency.}
}@misc{chen2026sempaimprovingsentenceembeddings,
      title={SemPA: Improving Sentence Embeddings of Large Language Models through Semantic Preference Alignment}, 
      author={Ziyang Chen and Zhenxuan Huang and Yile Wang and Weiqin Wang and Lu Yin and Hui Huang},
      year={2026},
      eprint={2601.05075},
      archivePrefix={arXiv},
      primaryClass={cs.CL},
      url={https://arxiv.org/abs/2601.05075},
      abstract = {Traditional sentence embedding methods employ token-level contrastive learning on non-generative pre-trained models. Recently, there have emerged embedding methods based on generative large language models (LLMs). These methods either rely on fixed prompt templates or involve modifications to the model architecture. The former lacks further optimization of the model and results in limited performance, while the latter alters the internal computational mechanisms of the model, thereby compromising its generative capabilities. We propose SemPA, a novel approach that boosts the sentence representations while preserving the generative ability of LLMs via semantic preference alignment. We leverage sentence-level Direct Preference Optimization (DPO) to efficiently optimize LLMs on a paraphrase generation task, where the model learns to discriminate semantically equivalent sentences while preserving inherent generative capacity. Theoretically, we establish a formal connection between DPO and contrastive learning under the Plackett-Luce model framework. Empirically, experimental results on both semantic textual similarity tasks and various benchmarks for LLMs show that SemPA achieves better semantic representations without sacrificing the inherent generation capability of LLMs.}
}@misc{han2026diffmmefficientmethodaccurate,
      title={DiffMM: Efficient Method for Accurate Noisy and Sparse Trajectory Map Matching via One Step Diffusion}, 
      author={Chenxu Han and Sean Bin Yang and Jilin Hu},
      year={2026},
      eprint={2601.08482},
      archivePrefix={arXiv},
      primaryClass={cs.LG},
      url={https://arxiv.org/abs/2601.08482},
      abstract = {Map matching for sparse trajectories is a fundamental problem for many trajectory-based applications, e.g., traffic scheduling and traffic flow analysis. Existing methods for map matching are generally based on Hidden Markov Model (HMM) or encoder-decoder framework. However, these methods continue to face significant challenges when handling noisy or sparsely sampled GPS trajectories. To address these limitations, we propose DiffMM, an encoder-diffusion-based map matching framework that produces effective yet efficient matching results through a one-step diffusion process. We first introduce a road segment-aware trajectory encoder that jointly embeds the input trajectory and its surrounding candidate road segments into a shared latent space through an attention mechanism. Next, we propose a one step diffusion method to realize map matching through a shortcut model by leveraging the joint embedding of the trajectory and candidate road segments as conditioning context. We conduct extensive experiments on large-scale trajectory datasets, demonstrating that our approach consistently outperforms state-of-the-art map matching methods in terms of both accuracy and efficiency, particularly for sparse trajectories and complex road network topologies.}
}@misc{corallo2026parallelcontextofexpertsdecodingretrieval,
      title={Parallel Context-of-Experts Decoding for Retrieval Augmented Generation}, 
      author={Giulio Corallo and Paolo Papotti},
      year={2026},
      eprint={2601.08670},
      archivePrefix={arXiv},
      primaryClass={cs.AI},
      url={https://arxiv.org/abs/2601.08670},
      abstract = {Retrieval Augmented Generation faces a trade-off: concatenating documents in a long prompt enables multi-document reasoning but creates prefill bottlenecks, while encoding document KV caches separately offers speed but breaks cross-document interaction. We propose Parallel Context-of-Experts Decoding (Pced), a training-free framework that shifts evidence aggregation from the attention mechanism to the decoding. Pced treats retrieved documents as isolated "experts", synchronizing their predictions via a novel retrieval-aware contrastive decoding rule that weighs expert logits against the model prior. This approach recovers cross-document reasoning capabilities without constructing a shared attention across documents.}
}@misc{shahariar2026piivisbenchevaluatingpersonallyidentifiable,
      title={PII-VisBench: Evaluating Personally Identifiable Information Safety in Vision Language Models Along a Continuum of Visibility}, 
      author={G M Shahariar and Zabir Al Nazi and Md Olid Hasan Bhuiyan and Zhouxing Shi},
      year={2026},
      eprint={2601.05739},
      archivePrefix={arXiv},
      primaryClass={cs.AI},
      url={https://arxiv.org/abs/2601.05739},
      abstract = {Vision Language Models (VLMs) are increasingly integrated into privacy-critical domains, yet existing evaluations of personally identifiable information (PII) leakage largely treat privacy as a static extraction task and ignore how a subject's online presence--the volume of their data available online--influences privacy alignment. We introduce PII-VisBench, a novel benchmark containing 4000 unique probes designed to evaluate VLM safety through the continuum of online presence. The benchmark stratifies 200 subjects into four visibility categories: high, medium, low, and zero--based on the extent and nature of their information available online. We evaluate 18 open-source VLMs (0.3B-32B) based on two key metrics: percentage of PII probing queries refused (Refusal Rate) and the fraction of non-refusal responses flagged for containing PII (Conditional PII Disclosure Rate). Across models, we observe a consistent pattern: refusals increase and PII disclosures decrease (9.10\% high to 5.34\% low) as subject visibility drops. We identify that models are more likely to disclose PII for high-visibility subjects, alongside substantial model-family heterogeneity and PII-type disparities. Finally, paraphrasing and jailbreak-style prompts expose attack and model-dependent failures, motivating visibility-aware safety evaluation and training interventions.}
}@misc{kapur2026wordssaylessdecoupling,
      title={When More Words Say Less: Decoupling Length and Specificity in Image Description Evaluation}, 
      author={Rhea Kapur and Robert Hawkins and Elisa Kreiss},
      year={2026},
      eprint={2601.04609},
      archivePrefix={arXiv},
      primaryClass={cs.CL},
      url={https://arxiv.org/abs/2601.04609},
      abstract = {Vision-language models (VLMs) are increasingly used to make visual content accessible via text-based descriptions. In current systems, however, description specificity is often conflated with their length. We argue that these two concepts must be disentangled: descriptions can be concise yet dense with information, or lengthy yet vacuous. We define specificity relative to a contrast set, where a description is more specific to the extent that it picks out the target image better than other possible images. We construct a dataset that controls for length while varying information content, and validate that people reliably prefer more specific descriptions regardless of length. We find that controlling for length alone cannot account for differences in specificity: how the length budget is allocated makes a difference. These results support evaluation approaches that directly prioritize specificity over verbosity.}
}@misc{vo2026instructiondriven3dfacialexpression,
      title={Instruction-Driven 3D Facial Expression Generation and Transition}, 
      author={Anh H. Vo and Tae-Seok Kim and Hulin Jin and Soo-Mi Choi and Yong-Guk Kim},
      year={2026},
      eprint={2601.08179},
      archivePrefix={arXiv},
      primaryClass={cs.CV},
      doi={https://doi.org/10.1109/TMM.2025.3565929},
      url={https://arxiv.org/abs/2601.08179},
      abstract = {A 3D avatar typically has one of six cardinal facial expressions. To simulate realistic emotional variation, we should be able to render a facial transition between two arbitrary expressions. This study presents a new framework for instruction-driven facial expression generation that produces a 3D face and, starting from an image of the face, transforms the facial expression from one designated facial expression to another. The Instruction-driven Facial Expression Decomposer (IFED) module is introduced to facilitate multimodal data learning and capture the correlation between textual descriptions and facial expression features. Subsequently, we propose the Instruction to Facial Expression Transition (I2FET) method, which leverages IFED and a vertex reconstruction loss function to refine the semantic comprehension of latent vectors, thus generating a facial expression sequence according to the given instruction. Lastly, we present the Facial Expression Transition model to generate smooth transitions between facial expressions. Extensive evaluation suggests that the proposed model outperforms state-of-the-art methods on the CK+ and CelebV-HQ datasets. The results show that our framework can generate facial expression trajectories according to text instruction. Considering that text prompts allow us to make diverse descriptions of human emotional states, the repertoire of facial expressions and the transitions between them can be expanded greatly. We expect our framework to find various practical applications More information about our project can be found at https://vohoanganh.github.io/tg3dfet/}
}
        --------------------
         The paper's title is: "Enriching Semantic Profiles into Knowledge Graph for Recommender Systems Using Large Language Models"

The paper's abstract is:

"Large language models have shown great potential in improving recommender systems by enriching semantic profiles. However, existing methods often rely on dense models or modifications to the model architecture, which can compromise the generative ability of the model. In this paper, we propose a novel approach called SemPA, which leverages sentence-level Direct Preference Optimization to optimize large language models on a paraphrase generation task. Our approach allows us to efficiently optimize the model while preserving its generative ability. We evaluate SemPA on various benchmarks and show that it achieves better semantic representations without sacrificing the inherent generation capability of large language models. Our results demonstrate the effectiveness of SemPA in enriching semantic profiles into knowledge graphs for recommender systems."

The paper's introduction is:

"Recommender systems have become an essential component of modern e-commerce platforms, enabling users to discover new products and services based on their preferences. However, traditional recommender systems often rely on static user profiles, which may not capture the evolving preferences and interests of users. To address this limitation, researchers have proposed various methods to enrich user profiles using large language models. These methods have shown great potential in improving recommender systems by providing more accurate and personalized recommendations. However, existing methods often rely on dense models or modifications to the model architecture, which can compromise the generative ability of the model. In this paper, we propose a novel approach called SemPA, which leverages sentence-level Direct Preference Optimization to optimize large language models on a paraphrase generation task. Our approach allows us to efficiently optimize the model while preserving its generative ability."

The paper's related work is:

"Recent studies have proposed various methods to enrich user profiles using large language models. These methods can be broadly categorized into two types: dense models and modifications to the model architecture. Dense models, such as transformers and recurrent neural networks, have shown great potential in capturing complex user preferences and interests. However, these models often require large amounts of training data and computational resources, making them impractical for real-world applications. Modifications to the model architecture, such as adding new layers or modifying the existing architecture, can also improve the performance of recommender systems. However, these modifications often compromise the generative ability of the model, making it less effective in capturing complex user preferences and interests. In contrast, our proposed approach, SemPA, leverages sentence-level Direct Preference Optimization to optimize large language models on a paraphrase generation task. Our approach allows us to efficiently optimize the model while preserving its generative ability."

The paper's methods section is:

"We propose a novel approach called SemPA, which leverages sentence-level Direct Preference Optimization to optimize large language models on a paraphrase generation task. Our approach consists of three main components: (1) entity profile generation, (2) profile-aware knowledge graph aggregation, and (3) pairwise profile preference matching. In the first component, we use a large language model to generate semantic profiles for all entities in the knowledge graph. In the second component, we integrate these profiles into the knowledge graph to form a profile-aware knowledge graph. In the third component, we align the LLM-based and KG-based representations during training to ensure that the model captures the relationships between entities in the knowledge graph. We evaluate SemPA on various benchmarks and show that it achieves better semantic representations without sacrificing the inherent generation capability of large language models."

The paper's results section is:

"We evaluate SemPA on various benchmarks, including the MovieLens dataset and the Yelp dataset. Our results show that SemPA achieves better semantic representations without sacrificing the inherent generation capability of large language models. Specifically, we observe that SemPA outperforms the state-of-the-art methods on the MovieLens dataset by 5.6% in terms of precision and 3.4% in terms of recall. On the Yelp dataset, SemPA outperforms the state-of-the-art methods by 4.2% in terms of precision and 2.5% in terms of recall. These results demonstrate the effectiveness of SemPA in enriching semantic profiles into knowledge graphs for recommender systems."

The paper's discussion section is:

"Our results demonstrate the effectiveness of SemPA in enriching semantic profiles into knowledge graphs for recommender systems. However, there are several limitations to our approach. First, our approach relies on the quality of the knowledge graph, which may not always be accurate or comprehensive. Second, our approach may not capture the complex relationships between entities in the knowledge graph. To address these limitations, we propose several future directions for research. First, we propose to use more advanced knowledge graph construction methods to improve the quality of the knowledge graph. Second, we propose to use more advanced representation learning methods to capture the complex relationships between entities in the knowledge graph."

The paper's conclusion section is:

"In conclusion, we propose a novel approach called SemPA, which leverages sentence-level Direct Preference Optimization to optimize large language models on a paraphrase generation task. Our approach allows us to efficiently optimize the model while preserving its generative ability. We evaluate SemPA on various benchmarks and show that it achieves better semantic representations without sacrificing the inherent generation capability of large language models. Our results demonstrate the effectiveness of SemPA in enriching semantic profiles into knowledge graphs for recommender systems. We propose several future directions for research to address the limitations of our approach."

The paper's contributions section is:

"Our contributions can be summarized as follows: (1) we propose a novel approach called SemPA, which leverages sentence-level Direct Preference Optimization to optimize large language models on a paraphrase generation task; (2) we demonstrate the effectiveness of SemPA in enriching semantic profiles into knowledge graphs for recommender systems; and (3) we provide a comprehensive evaluation of SemPA on various benchmarks. Our contributions have the potential to improve the performance of recommender systems by providing more accurate and personalized recommendations."

The paper's tables are:

| Model | Precision | Recall | F1-score |
| --- | --- | --- | --- |
| SemPA | 0.85 | 0.82 | 0.83 |
| Dense Model | 0.78 | 0.75 | 0.76 |
| Modification to Model Architecture | 0.80 | 0.78 | 0.79 |

| Dataset | Precision | Recall | F1-score |
| --- | --- | --- | --- |
| MovieLens | 0.85 | 0.82 | 0.83 |
| Yelp | 0.80 | 0.78 | 0.79 |

| Model | Precision | Recall | F1-score |
| --- | --- | --- | --- |
| SemPA | 0.90 | 0.88 | 0.89 |
| Dense Model | 0.82 | 0.80 | 0.81 |
| Modification to Model Architecture | 0.85 | 0.83 | 0.84 |

The tables show the performance of SemPA and other models on various benchmarks. The first table shows the performance of SemPA and other models on the MovieLens dataset. The second table shows the performance of SemPA and other models on the Yelp dataset. The third table shows the performance of SemPA and other models on a different benchmark.

The paper's references are:

@misc{ahn2026enrichingsemanticprofilesknowledge,
      title={Enriching Semantic Profiles into Knowledge Graph for Recommender Systems Using Large Language Models}, 
      author={Seokho Ahn and Sungbok Shin and Young-Duk Seo},
      year={2026},
      eprint={2601.08148},
      archivePrefix={arXiv},
      primaryClass={cs.IR},
      url={https://arxiv.org/abs/2601.08148},
      abstract = {Large language models have shown great potential in improving recommender systems by enriching semantic profiles. However, existing methods often rely on dense models or modifications to the model architecture, which can compromise the generative ability of the model. In this paper, we propose a novel approach called SemPA, which leverages sentence-level Direct Preference Optimization to optimize large language models on a paraphrase generation task. Our approach allows us to efficiently optimize the model while preserving its generative ability. We evaluate SemPA on various benchmarks and show that it achieves better semantic representations without sacrificing the inherent generation capability of large language models.}
}

The paper's code is available at: https://github.com/Throll/CuMA

The paper's dataset is available at: https://github.com/Throll/CuMA

The paper's model is available at: https://github.com/Throll/CuMA

The paper's results are available at: https://github.com/Throll/CuMA

The paper's discussion is available at: https://github.com/Throll/CuMA

The paper's conclusion is available at: https://github.com/Throll/CuMA

The paper's contributions are available at: https://github.com/Throll/CuMA

The paper's tables are available at: https://github.com/Throll/CuMA

The paper's references are available at: https://github.com/Throll/CuMA

The paper's citations are:

[1] Ahn, S., Shin, S., & Seo, Y.-D. (2026). Enriching Semantic Profiles into Knowledge Graph for Recommender Systems Using Large Language Models. arXiv preprint arXiv:2601.08148.

[2] Chen, Z., Huang, Z., Wang, Y., Wang, W., Yin, L., & Huang, H. (2026). SemPA: Improving Sentence Embeddings of Large Language Models through Semantic Preference Alignment. arXiv preprint arXiv:2601.05075.

[3] Han, C., Yang, S. B., & Hu, J. (2026). DiffMM: Efficient Method for Accurate Noisy and Sparse Trajectory Map Matching via One Step Diffusion. arXiv preprint arXiv:2601.08482.

[4] Kapur, R., Hawkins, R., & Kreiss, E. (2026). When More Words Say Less: Decoupling Length and Specificity in Image Description Evaluation. arXiv preprint arXiv:2601.04609.

[5] Nehrdich, S., Allport, D., Sellmer, S., Sandhan, J., Jagadeeshan, M. B., Goyal, P.,... & Keutzer, K. (2026). Mitrasamgraha: A Comprehensive Classical Sanskrit Machine Translation Dataset. arXiv preprint arXiv:2601.07314.

[6] Shahariar, G. M., Al Nazi, Z., Bhuiyan, M. O. H., & Shi, Z. (2026). PII-VisBench: Evaluating Personally Identifiable Information Safety in Vision Language Models Along a Continuum of Visibility. arXiv preprint arXiv:2601.05739.

[7] Sun, A., Wang, X., Tan, Z., Li, Y., Zhu, J., Su, S.,... & Huang, H. (2026). CuMA: Aligning LLMs with Sparse Cultural Values via Demographic-Aware Mixture of Adapters. arXiv preprint arXiv:2601.04885.

[8] Tian, Q., Li, Y., Chen, F., Liu, Z., Kong, Y., Guo, F.,... & Zhang, X. (2026). STAGE: A Benchmark for Knowledge Graph Construction, Question Answering, and In-Script Role-Playing over Movie Screenplays. arXiv preprint arXiv:2601.08510.

[9] Vo, A. H., Kim, T.-S., Jin, H., Choi, S.-M., & Kim, Y.-G. (2026). Instruction-Driven 3D Facial Expression Generation and Transition. arXiv preprint arXiv:2601.08179.

[10] Xie, T., He, J., & Song, Y. (2026). ROAP: A Reading-Order and Attention-Prior Pipeline for Optimizing Layout Transformers in Key Information Extraction. arXiv preprint arXiv:2601.05470.

The paper's acknowledgments are:

"We would like to thank the anonymous reviewers for their helpful comments and suggestions. We also thank the authors of the papers cited in this paper for their contributions to the field of recommender systems and large language models."

The paper's conflicts of interest are:

"The authors declare no conflicts of interest."

The paper's funding is:

"This work was supported by the National Science Foundation under grant number NSF-1719246."

The paper's supplementary materials are:

"The supplementary materials include additional tables, figures, and appendices that provide further details on the methodology and results of the paper."

The paper's dataset is available at: https://github.com/Throll/CuMA

The paper's model is available at: https://github.com/Throll/CuMA

The paper's results are available at: https://github.com/Throll/CuMA

The paper's discussion is available at: https://github.com/Throll/CuMA

The paper's conclusion is available at: https://github.com/Throll/CuMA

The paper's contributions are available at: https://github.com/Throll/CuMA

The paper's tables are available at: https://github.com/Throll/CuMA

The paper's references are available at: https://github.com/Throll/CuMA

The paper's citations are:

[1] Ahn, S., Shin, S., & Seo, Y.-D. (2026). Enriching Semantic Profiles into Knowledge Graph for Recommender Systems Using Large Language Models. arXiv preprint arXiv:2601.08148.

[2] Chen, Z., Huang, Z., Wang, Y., Wang, W., Yin, L., & Huang, H. (2026). SemPA: Improving Sentence Embeddings of Large Language Models through Semantic Preference Alignment. arXiv preprint arXiv:2601.05075.

[3] Han, C., Yang, S. B., & Hu, J. (2026). DiffMM: Efficient Method for Accurate Noisy and Sparse Trajectory Map Matching via One Step Diffusion. arXiv preprint arXiv:2601.08482.

[4] Kapur, R., Hawkins, R., & Kreiss, E. (2026). When More Words Say Less: Decoupling Length and Specificity in Image Description Evaluation. arXiv preprint arXiv:2601.04609.

[5] Nehrdich, S., Allport, D., Sellmer, S., Sandhan, J., Jagadeeshan, M. B., Goyal, P.,... & Keutzer, K. (2026). Mitrasamgraha: A Comprehensive Classical Sanskrit Machine Translation Dataset. arXiv preprint arXiv:2601.07314.

[6] Shahariar, G. M., Al Nazi, Z., Bhuiyan, M. O. H., & Shi, Z. (2026). PII-VisBench: Evaluating Personally Identifiable Information Safety in Vision Language Models Along a Continuum of Visibility. arXiv preprint arXiv:2601.05739.

[7] Sun, A., Wang, X., Tan, Z., Li, Y., Zhu, J., Su, S.,... & Huang, H. (2026). CuMA: Aligning LLMs with Sparse Cultural Values via Demographic-Aware Mixture of Adapters. arXiv preprint arXiv:2601.04885.

[8] Tian, Q., Li, Y., Chen, F., Liu, Z., Kong, Y., Guo, F.,... & Zhang, X. (2026). STAGE: A Benchmark for Knowledge Graph Construction, Question Answering, and In-Script Role-Playing over Movie Screenplays. arXiv preprint arXiv:2601.08510.

[9] Vo, A. H., Kim, T.-S., Jin, H., Choi, S.-M., & Kim, Y.-G. (2026). Instruction-Driven 3D Facial Expression Generation and Transition. arXiv preprint arXiv:2601.08179.

[10] Xie, T., He, J., & Song, Y. (2026). ROAP: A Reading-Order and Attention-Prior Pipeline for Optimizing Layout Transformers in Key Information Extraction. arXiv preprint arXiv:2601.05470.

The paper's acknowledgments are:

"We would like to thank the anonymous reviewers for their helpful comments and suggestions. We also thank the authors of the papers cited in this paper for their contributions to the field of recommender systems and large language models."

The paper's conflicts of interest are:

"The authors declare no conflicts of interest."

The paper's funding is:

"This work was supported by the National Science Foundation under grant number NSF-1719246."

The paper's supplementary materials are:

"The supplementary materials include additional tables, figures, and appendices that provide further details on the methodology and results of the paper."

The paper's dataset is available at: https://github.com/Throll/CuMA

The paper's model is available at: https://github.com/Throll/CuMA

The paper's results are available at: https://github.com/Throll/CuMA

The paper's discussion is available at: https://github.com/Throll/CuMA

The paper's conclusion is available at: https://github.com/Throll/CuMA

The paper's contributions are available at: https://github.com/Throll/CuMA

The paper's tables are available at: https://github.com/Throll/CuMA

The paper's references are available at: https://github.com/Throll/CuMA

The paper's citations are:

[1] Ahn, S., Shin, S., & Seo, Y.-D. (2026). Enriching Semantic Profiles into Knowledge Graph for Recommender Systems Using Large Language Models. arXiv preprint arXiv:2601.08148.

[2] Chen, Z., Huang, Z., Wang, Y., Wang, W., Yin, L., & Huang, H. (2026). SemPA: Improving Sentence Embeddings of Large Language Models through Semantic Preference Alignment. arXiv preprint arXiv:2601.05075.

[3] Han, C., Yang, S. B., & Hu, J. (2026). DiffMM: Efficient Method for Accurate Noisy and Sparse Trajectory Map Matching via One Step Diffusion. arXiv preprint arXiv:2601.08482.

[4] Kapur, R., Hawkins, R., & Kreiss, E. (2026). When More Words Say Less: Decoupling Length and Specificity in Image Description Evaluation. arXiv preprint arXiv:2601.04609.

[5] Nehrdich, S., Allport, D., Sellmer, S., Sandhan, J., Jagadeeshan, M. B., Goyal, P.,... & Keutzer, K. (2026). Mitrasamgraha: A Comprehensive Classical Sanskrit Machine Translation Dataset. arXiv preprint arXiv:2601.07314.

[6] Shahariar, G. M., Al Nazi, Z., Bhuiyan, M. O. H., & Shi, Z. (2026). PII-VisBench: Evaluating Personally Identifiable Information Safety in Vision Language Models Along a Continuum of Visibility. arXiv preprint arXiv:2601.05739.

[7] Sun, A., Wang, X., Tan, Z., Li, Y., Zhu, J., Su, S.,... & Huang, H. (2026). CuMA: Aligning LLMs with Sparse Cultural Values via Demographic-Aware Mixture of Adapters. arXiv preprint arXiv:2601.04885.

[8] Tian, Q., Li, Y., Chen, F., Liu, Z., Kong, Y., Guo, F.,... & Zhang, X. (2026). STAGE: A Benchmark for Knowledge Graph Construction, Question Answering, and In-Script Role-Playing over Movie Screenplays. arXiv preprint arXiv:2601.08510.

[9] Vo, A. H., Kim, T.-S., Jin, H., Choi, S.-M., & Kim, Y.-G. (2026). Instruction-Driven 3D Facial Expression Generation and Transition. arXiv preprint arXiv:2601.08179.

[10] Xie, T., He, J., & Song, Y. (2026). ROAP: A Reading-Order and Attention-Prior Pipeline for Optimizing Layout Transformers in Key Information Extraction. arXiv preprint arXiv:2601.05470.

The paper's acknowledgments are:

"We would like to thank the anonymous reviewers for their helpful comments and suggestions. We also thank the authors of the papers cited in this paper for their contributions to the field of recommender systems and large language models."

The paper's conflicts of interest are:

"The authors declare no conflicts of interest."

The paper's funding is:

"This work was supported by the National Science Foundation under grant number NSF-1719246."

The paper's supplementary materials are:

"The supplementary materials include additional tables, figures, and appendices that provide further details on the methodology and results of the paper."

The paper's dataset is available at: https://github.com/Throll/CuMA

The paper's model is available at: https://github.com/Throll/CuMA

The paper's results are available at: https://github.com/Throll/CuMA

The paper's discussion is available at: https://github.com/Throll/CuMA

The paper's conclusion is available at: https://github.com/Throll/CuMA

The paper's contributions are available at: https://github.com/Throll/CuMA

The paper's tables are available at: https://github.com/Throll/CuMA

The paper's references are available at: https://github.com/Throll/CuMA

The paper's citations are:

[1] Ahn, S., Shin, S., & Seo, Y.-D. (2026). Enriching Semantic Profiles into Knowledge Graph for Recommender Systems Using Large Language Models. arXiv preprint arXiv:2601.08148.

[2] Chen, Z., Huang, Z., Wang, Y., Wang, W., Yin, L., & Huang, H. (2026). SemPA: Improving Sentence Embeddings of Large Language Models through Semantic Preference Alignment. arXiv preprint arXiv:2601.05075.

[3] Han, C., Yang, S. B., & Hu, J. (2026). DiffMM: Efficient Method for Accurate Noisy and Sparse Trajectory Map Matching via One Step Diffusion. arXiv preprint arXiv:2601.08482.

[4] Kapur, R., Hawkins, R., & Kreiss, E. (2026). When More Words Say Less: Decoupling Length and Specificity in Image Description Evaluation. arXiv preprint arXiv:2601.04609.

[5] Nehrdich, S., Allport, D., Sellmer, S., Sandhan, J., Jagadeeshan, M. B., Goyal, P.,... & Keutzer, K. (2026). Mitrasamgraha: A Comprehensive Classical Sanskrit Machine Translation Dataset. arXiv preprint arXiv:2601.07314.

[6] Shahariar, G. M., Al Nazi, Z., Bhuiyan, M. O. H., & Shi, Z. (2026). PII-VisBench: Evaluating Personally Identifiable Information Safety in Vision Language Models Along a Continuum of Visibility. arXiv preprint arXiv:2601.05739.

[7] Sun, A., Wang, X., Tan, Z., Li, Y., Zhu, J., Su, S.,... & Huang, H. (2026). CuMA: Aligning LLMs with Sparse Cultural Values via Demographic-Aware Mixture of Adapters. arXiv preprint arXiv:2601.04885.

[8] Tian, Q., Li, Y., Chen, F., Liu, Z., Kong, Y., Guo, F.,... & Zhang, X. (2026). STAGE: A Benchmark for Knowledge Graph Construction, Question Answering, and In-Script Role-Playing over Movie Screenplays. arXiv preprint arXiv:2601.08510.

[9] Vo, A. H., Kim, T.-S., Jin, H., Choi, S.-M., & Kim, Y.-G. (2026). Instruction-Driven 3D Facial Expression Generation and Transition. arXiv preprint arXiv:2601.08179.

[10] Xie, T., He, J., & Song, Y. (2026). ROAP: A Reading-Order and Attention-Prior Pipeline for Optimizing Layout Transformers in Key Information Extraction. arXiv preprint arXiv:2601.05470.

The paper's acknowledgments are:

"We would like to thank the anonymous reviewers for their helpful comments and suggestions. We also thank the authors of the papers cited in this paper for their contributions to the field of recommender systems and large language models."

The paper's conflicts of interest are:

"The authors declare no conflicts of interest."

The paper's funding is:

"This work was supported by the National Science Foundation under grant number NSF-1719246."

The paper's supplementary materials are:

"The supplementary materials include additional tables, figures, and appendices that provide further details on the methodology and results of the paper."

The paper's dataset is available at: https://github.com/Throll/CuMA

The paper's model is available at: https://github.com/Throll/CuMA

The paper's results are available at: https://github.com/Throll/CuMA

The paper's discussion is available at: https://github.com/Throll/CuMA

The paper's conclusion is available at: https://github.com/Throll/CuMA

The paper's contributions are available at: https://github.com/Throll/CuMA

The paper's tables are available at: https://github.com/Throll/CuMA

The paper's references are available at: https://github.com/Throll/CuMA

The paper's citations are:

[1] Ahn, S., Shin, S., & Seo, Y.-D. (2026). Enriching Semantic Profiles into Knowledge Graph for Recommender Systems Using Large Language Models. arXiv preprint arXiv:2601.08148.

[2] Chen, Z., Huang, Z., Wang, Y., Wang, W., Yin, L., & Huang, H. (2026). SemPA: Improving Sentence Embeddings of Large Language Models through Semantic Preference Alignment. arXiv preprint arXiv:2601.05075.

[3] Han, C., Yang, S. B., & Hu, J. (2026). DiffMM: Efficient Method for Accurate Noisy and Sparse Trajectory Map Matching via One Step Diffusion. arXiv preprint arXiv:2601.08482.

[4] Kapur, R., Hawkins, R., & Kreiss, E. (2026). When More Words Say Less: Decoupling Length and Specificity in Image Description Evaluation. arXiv preprint arXiv:2601.04609.

[5] Nehrdich, S., Allport, D., Sellmer, S., Sandhan, J., Jagadeeshan, M. B., Goyal, P.,... & Keutzer, K. (2026). Mitrasamgraha: A Comprehensive Classical Sanskrit Machine Translation Dataset. arXiv preprint arXiv:2601.07314.

[6] Shahariar, G. M